\section{The Amorphous Kitaev Model}

In Chapter \ref{chap:amorphous_kitaev} we turned our attention to the Kitaev honeycomb model \cite{kitaev_anyons_2006}. We started by providing a review of the model, and explaining the Majorana fermion construction used to find an exact solution, showing how the Hamiltonian may be re-expressed in terms of a $\mathbb Z_2$ gauge field, and discussed the role of time reversal symmetry. We presented the full ground-state phase diagram containing the three `toric code' $A_{\alpha}$ phases, as well as the gapless $B$ phase and discussed how the addition of a magnetic field can serve to open a gap, allowing the phase to host non-Abelian anyonic excitations. In \textsection\ref{sec:amorphous_kitaev_model} we presented a general method for constructing and solving the Kitaev Hamiltonian, using the Majorana method, on an arbitrary three-coordinated lattice. This included the method for assigning a valid edge-colouring, as well as showing how the $\mathbb Z_2$ fluxes can be calculated on an arbitrary lattice and providing numerical evidence for the conjecture that the ground state flux sector could be found in any lattice according to a simple counting rule. We presented the full ground-state phase diagram and showed that an analogous phase to Kitaev's original non-Abelian $B$ phase exists. We calculated the Chern number using the local markers discussed in the previous chapter, and demonstrated the presence of edge modes in this phase. Finally, we discussed the role of disordering the coupling strengths, showing that the gapped phases were robust against such disorder. Finally, we concluded that the Kitaev spin liquid ground state is robust on a general lattice, providing the first example of an amorphous quantum spin liquid.

Of course, the most obvious next step is to look for a physical material that might be able to realise a Hamiltonian with dominant Kitaev-type interactions on an amorphous lattice. There are a number of promising candidates, for example it has recently been shown that the Kitaev material Li\textsubscript{2}IrO\textsubscript{3} can be created with amorphous structure \cite{lee_lithiair_2019}. In fact, by using the standard method of repeated liquifying and fast cooling any crystalline material can be made amorphous \cite{zallen_physics_2008,turnbull_under_1969, kaneyoshi_amorphous_2018,petrakovskii_amorphous_1981}, opening the possibility for exploring a wide variety of non-crystalline Kitaev materials. We expect that an experimental signature of a chiral amorphous QSL would be a half-quantized thermal Hall effect similar to the magnetic field induced behaviour of honeycomb Kitaev materials~\cite{kasahara_majorana_2018,yokoi_majorana_2021,yamashita_sample_2020,bruin_robustness_2022}.

Every result presented in \textsection\ref{sec:amorphous_kitaev_model} depends on our ansatz for the ground-state flux sector, which is an extension of the theorem proved by Lieb in 1994 \cite{lieb_flux_1994}. However, this theorem remains unproved, and we were only able to collect compelling numerical evidence to support the conjecture. In \textsection\ref{sec:cycles_and_ground_state} we provided some analytical progress, suggesting that -- at least at the lowest order that each plaquette appears in the energy -- Lieb's ansatz seems correct. Thus, a very interesting next step would be to find a route towards proving this theorem. 

In our analysis, we largely focussed on the non-Abelian $B$ phase, however there is good reason to suspect that new physics could be found also by considering the `toric code' phases found by setting only one bond type to be much stronger than the other two. This would effectively form a dimerisation of the lattice, with the dimers weakly interacting through the remaining two bond types. Kitaev showed in the original work that these weak interactions could be treated using perturbation theory, arriving at an effective Hamiltonian where the couplings formed loop operators around each (square -- since two bonds in each hexagon formed a dimer) plaquette. This provided a mapping onto the toric code model \cite{kitaev_fault-tolerant_2003}. The Kitaev model has been studied extensively in this limit, on the honeycomb as well as star lattice \cite{dusuel_creation_2008,schmidt_emergent_2008,vidal_perturbative_2008}, however no work has been done to understand the amorphous case. This is particularly exciting since the order in perturbation that a loop operator appears depends on the number of sides in the constituent plaquette, meaning that the amorphous lattice has the potential to realise a much richer and more complex effective system than the honeycomb case. 

Finally, we also note that the classical Kitaev model warrants investigation in an amorphous context. Firstly, the degenerate manifold of classical states presented in \textsection\ref{sec:classical_kitaev} can be straightforwardly extended to the amorphous model. In the honeycomb case, even the $S \gg 1$ model has been shown to form a quantum spin liquid state -- remarkably also equivalent to the toric code \cite{baskaran_spin_2008, rousochatzakis_quantum_2018}. The construction of this QSL phase relied extensively on the symmetries of the lattice: the system preferentially orders in a small subset of the Cartesian states due to the effect of fluctuations, forming a honeycomb super-lattice of frozen spin dimers. Each of these pseudospins then act as an effective spin-$\frac 12$ interacting according to a toric code Hamiltonian. Due to the strong dependence of this argument on the symmetries of the honeycomb lattice, it would be worth investigating which states are selected by fluctuations on an amorphous lattice, and whether these models are also able to realise an effective toric code phase. 